\documentclass{../../../unipd-notes}

\unipdsetup{
  course = {Fondamenti di Ingegneria Informatica},
  short-course = {Fondamenti di Ingegneria},
  document-type = {Appunti delle lezioni}
}

\begin{document}
\makecoursetableofcontents

\chapter{Architettura del calcolatore}
\section{Unità centrale di elaborazione}
La CPU esegue le istruzioni attraverso le fasi di prelievo, decodifica ed esecuzione.
\subsection{Registri}
I registri conservano operandi, indirizzi e stato della computazione corrente.
\section{Gerarchia di memoria}
Cache, memoria principale e memoria secondaria bilanciano capacità e latenza.

\chapter{Sistemi operativi}
\section{Processi e concorrenza}
Il sistema operativo assegna la CPU ai processi pronti mediante lo scheduler.
\end{document}
